\documentclass[12pt,a4paper]{article}
% Drake_preamble.tex - LaTeX Preamble Template
% 作者: 謝慕揚, MD, PhD, FESC
% 
% 使用說明:
% 1. 表格請使用 longtable，不要有垂直分隔線 |
% 2. 表格內換行使用 \newline，不要使用 \\
% 3. 藥物名稱、檢驗、檢查請維持英文
% 4. Reference 格式請使用 \href 加上驗證過的 DOI

% Including essential packages for Chinese typesetting
\usepackage{xeCJK}
\usepackage{fontspec}
% Configuring Chinese font with Noto Serif CJK TC, fallback to SimSun
\IfFontExistsTF{Noto Serif CJK TC}{
  \setCJKmainfont{Noto Serif CJK TC}
  \setCJKsansfont{Noto Serif CJK TC}
  \setCJKmonofont{Noto Serif CJK TC}
}{\setCJKmainfont{SimSun}
  \setCJKsansfont{SimSun}
  \setCJKmonofont{SimSun}
}

% Including other necessary packages
\usepackage{geometry}
\usepackage{titlesec}
\usepackage{enumitem}
\usepackage{xcolor}
\usepackage{colortbl}
\usepackage{tcolorbox}
\usepackage{booktabs}
\usepackage{longtable}
\usepackage{array}
\usepackage{graphicx}
\usepackage{tikz}
\usepackage{fancyhdr}
\usepackage{multicol}
\usepackage{datetime2}
\usepackage{hyperref}
\usepackage{mdframed}
\usepackage{amsmath}
\usepackage{amssymb}
\usepackage{multirow}

\hypersetup{
    colorlinks=true,
    linkcolor=emergency_blue,
    citecolor=emergency_blue,
    urlcolor=emergency_blue,
    pdfborder={0 0 0}
}

% Defining page geometry
\geometry{
  top=2.5cm,
  bottom=2.5cm,
  left=2cm,
  right=2cm,
  headheight=25pt
}

% Adjusting topmargin to compensate for increased headheight
\addtolength{\topmargin}{-10pt}

% Defining colors
\definecolor{emergency_red}{RGB}{186,24,27}
\definecolor{emergency_blue}{RGB}{0,114,188}
\definecolor{emergency_gray}{RGB}{120,120,120}
\definecolor{highlight_yellow}{RGB}{255,250,205}

% Setting title formats
\titleformat{\section}[block]{\Large\bfseries\color{emergency_red}}{\thesection}{10pt}{}
\titleformat{\subsection}[block]{\large\bfseries\color{emergency_blue}}{\thesubsection}{8pt}{}
\titleformat{\subsubsection}[block]{\normalsize\bfseries\color{emergency_gray}}{\thesubsubsection}{6pt}{}

% Defining custom environment for important notes
\newmdenv[
    linecolor=emergency_red,
    backgroundcolor=highlight_yellow,
    linewidth=2pt,
    topline=true,
    bottomline=true,
    leftline=true,
    rightline=true,
    innertopmargin=10pt,
    innerbottommargin=10pt,
    innerleftmargin=10pt,
    innerrightmargin=10pt
]{importantbox}

% Defining note command with tcolorbox
\newcommand{\note}[1]{
    \par
    \noindent
    \begin{tcolorbox}[
        colback=emergency_blue!5!white,
        colframe=emergency_blue,
        title=\textbf{重點提醒},
        fonttitle=\bfseries,
        boxrule=0.5pt,
        arc=3pt,
        left=6pt,
        right=6pt,
        top=6pt,
        bottom=6pt
    ]
    #1
    \end{tcolorbox}
}

% Key findings box
\newcommand{\keyfinding}[1]{
    \par
    \noindent
    \begin{tcolorbox}[
        colback=yellow!10!white,
        colframe=orange!75!black,
        title=\textbf{核心發現摘要},
        fonttitle=\bfseries,
        boxrule=0.5pt,
        arc=3pt
    ]
    #1
    \end{tcolorbox}
}

% Defining custom date format for ISO 8601
\DTMnewdatestyle{iso}{
  \renewcommand{\DTMdisplaydate}[4]{\number##1-\number##2-\number##3}
  \renewcommand{\DTMDisplaydate}[4]{\number##1-\number##2-\number##3}
}
\DTMsetdatestyle{iso}
\newcommand{\mydate}{\DTMtoday}


% Custom header for this document
\pagestyle{fancy}
\fancyhf{}
\fancyhead[L]{\textcolor{emergency_red}{Intensive Care Medicine 教學講義}}
\fancyhead[R]{\textcolor{emergency_blue}{Stress Ulcer Prevention}}
\fancyfoot[C]{\textcolor{emergency_gray}{\thepage}}
\renewcommand{\headrulewidth}{1pt}
\renewcommand{\headrule}{\hbox to\headwidth{\color{emergency_red}\leaders\hrule height \headrulewidth\hfill}}

\begin{document}

\begin{center}
{\LARGE\bfseries\textcolor{emergency_red}{Intensive Care Medicine Editorial}}\\[0.5cm]
{\Large\textbf{Preventing Stress Ulcer Bleeding}}\\[0.3cm]
{\large 預防壓力性潰瘍出血}\\[0.5cm]
{\normalsize Intensive Care Med 2024;50:2162-2165. DOI: 10.1007/s00134-024-07674-3}\\[0.3cm]
{\normalsize 整理：謝慕揚 MD, PhD, FESC \quad 日期：\mydate}
\end{center}

\keyfinding{
\textbf{核心訊息：}Proton pump inhibitors (PPIs) 可有效減少 clinically important 及 patient-important upper GI bleeding，但對 mortality 的影響不確定。

\textbf{臨床建議：}
\begin{itemize}
    \item \textbf{High certainty evidence：}PPIs 可減少上消化道出血
    \item \textbf{需注意：}High illness acuity 病人使用 PPIs 可能與 increased mortality 相關
    \item \textbf{停藥時機：}ICU 出院後應積極評估停藥，permanent prescription 無 trial data 支持
\end{itemize}
}

\tableofcontents
\newpage

%==============================================================
\section{文獻概述}
%==============================================================

本篇為 Intensive Care Medicine 的 Editorial，由紐西蘭 Wellington Hospital 的 Young 等人撰寫，整合 REVISE 及 SUP-ICU 等大型試驗結果，討論 stress ulcer prophylaxis 的實務應用。

\subsection{作者資訊}
\textbf{通訊作者：}Paul J. Young, MD\\
\textbf{機構：}Intensive Care Unit, Wellington Hospital, Wellington, New Zealand; Australian and New Zealand Intensive Care Research Centre, Monash University, Melbourne, Australia

%==============================================================
\section{背景知識}
%==============================================================

\subsection{Stress Ulcers 定義與影響}

\begin{itemize}
    \item \textbf{定義：}Gastric 或 duodenal mucosal erosions，是 critical illness 的已知併發症
    \item \textbf{表現範圍：}
    \begin{itemize}
        \item Minor bleeding（如暫時性 coffee ground gastric aspirates）
        \item Life-threatening gastrointestinal bleeding
    \end{itemize}
    \item \textbf{臨床後果：}
    \begin{itemize}
        \item Patient and family concerns
        \item 需要檢查（如 endoscopy）
        \item 需要治療（如 blood transfusions）
        \item 相關 morbidity 及 mortality
    \end{itemize}
\end{itemize}

\subsection{重要證據摘要}

\note{
\textbf{關鍵發現：}
\begin{itemize}
    \item \textbf{High certainty evidence：}PPIs 可有效減少 clinically important 及 patient-important upper GI bleeding
    \item \textbf{Uncertain：}Stress ulcer prophylaxis 對 mortality 的影響
    \item \textbf{Rare：}直接歸因於 stress ulceration GI bleeding 的死亡
\end{itemize}
}

%==============================================================
\section{藥物選擇}
%==============================================================

\begin{longtable}{p{4cm}p{10cm}}
\toprule
\textbf{藥物類別} & \textbf{說明} \\
\midrule
\endfirsthead

\toprule
\textbf{藥物類別} & \textbf{說明} \\
\midrule
\endhead

\bottomrule
\endfoot

\textbf{Proton Pump Inhibitors (PPIs)} &
• Stress ulcer prophylaxis 的主要藥物\newline
• 效果優於 H2-receptor blockers \\
\midrule
\textbf{H2-Receptor Blockers} &
• 較少使用\newline
• 效果不如 PPIs\newline
• Ranitidine（此適應症研究最多的藥物）已下市 \\
\bottomrule
\end{longtable}

%==============================================================
\section{處方策略}
%==============================================================

\subsection{目前實務現況}

\begin{itemize}
    \item Stress ulcer prophylaxis 已納入：
    \begin{itemize}
        \item ICU admission order sets
        \item Checklists（如「FAST HUG」）
        \item Guidelines
    \end{itemize}
    \item 許多單位採取近乎 \textbf{universal administration} 的做法
\end{itemize}

\subsection{處方方式比較}

\begin{longtable}{p{4cm}p{10cm}}
\toprule
\textbf{方式} & \textbf{說明} \\
\midrule
\endfirsthead

\toprule
\textbf{方式} & \textbf{說明} \\
\midrule
\endhead

\bottomrule
\endfoot

\textbf{Universal Administration} &
• 幾乎所有 ICU 病人都給予\newline
• 簡單執行 \\
\midrule
\textbf{Individualized Approach} &
• 每日評估 bleeding risks vs prophylaxis benefits\newline
• 考量 pre-existing 及 dynamic bleeding risks \\
\midrule
\textbf{Risk-Based Approach 的限制} &
• GI bleeding 常發生在無 identifiable risk factors 的病人\newline
• 無法預測 bleeding episodes 的嚴重程度 \\
\bottomrule
\end{longtable}

%==============================================================
\section{臨床試驗納入條件}
%==============================================================

\subsection{SUP-ICU Trial 納入條件}

所有成人（≥18 歲）急性入住 ICU 且具有以下一項或多項 GI bleeding 風險因子：

\begin{longtable}{p{4cm}p{10cm}}
\toprule
\textbf{風險因子} & \textbf{定義} \\
\midrule
\endfirsthead

\toprule
\textbf{風險因子} & \textbf{定義} \\
\midrule
\endhead

\bottomrule
\endfoot

\textbf{Shock} &
• 持續使用 vasopressors/inotropes\newline
• SBP <90 mmHg\newline
• MAP <70 mmHg\newline
• Lactate ≥4 mmol/L \\
\midrule
\textbf{RRT} & Acute 或 chronic intermittent 或 CRRT \\
\midrule
\textbf{Mechanical Ventilation} & Invasive MV 預期 >24 小時 \\
\midrule
\textbf{Coagulopathy} &
• Platelets <50×10$^9$/L\newline
• INR >1.5\newline
• PT >20 秒\newline
（24 小時內記錄）\\
\midrule
\textbf{Anticoagulant Treatment} & 進行中的 anticoagulant 治療（prophylactic doses 除外）\\
\midrule
\textbf{Coagulopathy History} & 入院前 6 個月內有 coagulopathy 病史 \\
\midrule
\textbf{Chronic Liver Disease} &
• Portal hypertension\newline
• Cirrhosis（biopsy/CT/US 證實）\newline
• Variceal bleeding 病史\newline
• Hepatic encephalopathy 病史 \\
\bottomrule
\end{longtable}

\subsection{REVISE Trial 納入條件}

\begin{itemize}
    \item 成人 ≥18 歲
    \item 接受 invasive mechanical ventilation
    \item 預期 mechanical ventilation 將持續超過 randomization 後的隔天
\end{itemize}

\subsection{排除 PPI Prophylaxis 的情況}

\begin{itemize}
    \item PPI 或其添加物過敏史
    \item 有 active GI bleeding 或 high bleeding risk：
    \begin{itemize}
        \item Current bleeding
        \item 8 週內 peptic ulcer bleeding
        \item Recent severe esophagitis
        \item Barrett's esophagus
        \item Zollinger-Ellison syndrome
    \end{itemize}
    \item 與 PPI 有交互作用的特定藥物（如 HIV protease inhibitors atazanavir 或 nelfinavir）
\end{itemize}

%==============================================================
\section{Heterogeneity of Treatment Effect}
%==============================================================

\note{
\textbf{重要發現：}PPI 對 mortality 的影響可能因 baseline illness severity 而異！
}

\subsection{SUP-ICU Trial 次群體分析}

\begin{longtable}{p{5cm}p{9cm}}
\toprule
\textbf{Illness Acuity} & \textbf{90-day Mortality RR (PPI vs Placebo)} \\
\midrule
\endfirsthead

\toprule
\textbf{Illness Acuity} & \textbf{90-day Mortality RR (PPI vs Placebo)} \\
\midrule
\endhead

\bottomrule
\endfoot

Lower illness acuity & 0.92 (95\% CI, 0.78-1.09) \\
\midrule
Higher illness acuity & 1.13 (95\% CI, 0.99-1.30) \\
\midrule
\textbf{Interaction P} & \textbf{0.05} \\
\bottomrule
\end{longtable}

\subsection{REVISE Trial 次群體分析}

\begin{longtable}{p{5cm}p{9cm}}
\toprule
\textbf{Illness Acuity} & \textbf{90-day Mortality HR (PPI vs Placebo)} \\
\midrule
\endfirsthead

\toprule
\textbf{Illness Acuity} & \textbf{90-day Mortality HR (PPI vs Placebo)} \\
\midrule
\endhead

\bottomrule
\endfoot

Lower illness acuity & 0.85 (95\% CI, 0.74-0.98) \\
\midrule
Higher illness acuity & 1.04 (95\% CI, 0.89-1.20) \\
\midrule
\textbf{Interaction P} & \textbf{0.27} \\
\bottomrule
\end{longtable}

\subsection{PEPTIC Trial 發現}

\begin{itemize}
    \item 比較 PPIs vs H2-receptor blockers
    \item Allocation to PPIs 在 \textbf{high illness acuity} 病人與 \textbf{increased mortality} 相關
    \item PPIs 如何增加 severely ill patients mortality 的機轉仍為推測
\end{itemize}

\note{
\textbf{臨床矛盾：}
\begin{itemize}
    \item PPIs 預防 upper GI bleeding 的效果在 high illness acuity 病人最顯著
    \item 但同樣是 high illness acuity 病人，PPIs 可能與 increased mortality 相關
\end{itemize}
}

%==============================================================
\section{Enteral Nutrition 的角色}
%==============================================================

\begin{itemize}
    \item Enteral nutrition 是否能降低 bleeding risk：\textbf{不確定}
    \item 在兩個大多數病人接受 enteral nutrition 的當代試驗中，PPIs 仍顯著減少 GI bleeding
\end{itemize}

\subsection{SUP-ICU Trial Post Hoc Analysis}

\begin{longtable}{p{5cm}p{9cm}}
\toprule
\textbf{項目} & \textbf{發現} \\
\midrule
\endfirsthead

\toprule
\textbf{項目} & \textbf{發現} \\
\midrule
\endhead

\bottomrule
\endfoot

Enteral nutrition 的效果 &
• 減少 GI bleeding\newline
• 增加 pneumonia\newline
• 減少 mortality \\
\midrule
Pantoprazole + Enteral nutrition &
• 可能與 increased all-cause mortality 相關\newline
• HR 1.27 (95\% CI, 0.99-1.64), P=0.06\newline
• Interaction P=0.02 \\
\bottomrule
\end{longtable}

\begin{itemize}
    \item 目前 data 不足以確認 enteral nutrition 建立後即可停止 stress ulcer prophylaxis
    \item 需要更多研究
\end{itemize}

%==============================================================
\section{停藥時機}
%==============================================================

\note{
\textbf{重要提醒：}
\begin{itemize}
    \item ICU 出院後無明確適應症的 \textbf{permanent prescription} 無 trial data 支持
    \item 一項大型 propensity-matched study 顯示，這類病人有較高風險：
    \begin{itemize}
        \item Pneumonia
        \item Cardiovascular events
        \item Hospital readmission
        \item 2-year mortality
    \end{itemize}
\end{itemize}
}

\subsection{建議}

\begin{itemize}
    \item 如同所有藥物，prophylaxis 應 \textbf{每日評估}
    \item 特別是在 \textbf{transitions in care} 時：
    \begin{itemize}
        \item 轉至一般病房時
        \item 出院時
    \end{itemize}
\end{itemize}

%==============================================================
\section{未來發展}
%==============================================================

\begin{itemize}
    \item 應用 \textbf{machine learning models} 分析 REVISE 及 SUP-ICU 等 RCT datasets
    \item 可能提供 \textbf{individualized treatment effects} 的估計
    \item 預測：
    \begin{itemize}
        \item 哪些病人最可能從 stress ulcer prophylaxis 獲益（預防 upper GI bleed）
        \item 哪些病人最不可能因 prophylaxis 而增加死亡風險
    \end{itemize}
    \item \textbf{Individualized treatment decisions} 可能即將實現
\end{itemize}

%==============================================================
\section{Take-Home Message}
%==============================================================

\begin{enumerate}
    \item \textcolor{red}{\textbf{有效預防出血：}}PPIs 可有效減少 clinically important 及 patient-important upper GI bleeding（high certainty evidence）
    \item \textcolor{red}{\textbf{Mortality 效果不確定：}}Stress ulcer prophylaxis 對 mortality 的影響不確定，可能因 illness severity 而異
    \item \textcolor{red}{\textbf{High acuity 病人需注意：}}High illness acuity 病人使用 PPIs 可能與 increased mortality 相關
    \item \textcolor{red}{\textbf{積極評估停藥：}}ICU 出院後應評估停藥，permanent prescription 無 evidence 支持且可能有害
\end{enumerate}

%==============================================================
\section{縮寫對照表}
%==============================================================

\begin{longtable}{p{3cm}p{11cm}}
\toprule
\textbf{縮寫} & \textbf{全名} \\
\midrule
\endfirsthead

\toprule
\textbf{縮寫} & \textbf{全名} \\
\midrule
\endhead

\bottomrule
\endfoot

PPI & Proton Pump Inhibitor (質子幫浦抑制劑) \\
GI & Gastrointestinal (胃腸道) \\
ICU & Intensive Care Unit (加護病房) \\
RRT & Renal Replacement Therapy (腎臟替代療法) \\
CRRT & Continuous Renal Replacement Therapy \\
MV & Mechanical Ventilation (機械通氣) \\
INR & International Normalized Ratio \\
PT & Prothrombin Time \\
SBP & Systolic Blood Pressure \\
MAP & Mean Arterial Pressure \\
HR & Hazard Ratio \\
RR & Relative Risk \\
CI & Confidence Interval \\
\bottomrule
\end{longtable}

%==============================================================
\section{參考文獻}
%==============================================================

\begin{enumerate}
    \item Young PJ, Cook DJ, Deane AM. Preventing stress ulcer bleeding. \href{https://doi.org/10.1007/s00134-024-07674-3}{\textit{Intensive Care Med}. 2024;50:2162-2165.}

    \item Cook D, Deane A, Lauzier F, et al. Stress ulcer prophylaxis during invasive mechanical ventilation. \href{https://doi.org/10.1056/NEJMoa2404245}{\textit{N Engl J Med}. 2024;391:9-20.}

    \item Krag M, Marker S, Perner A, et al. Pantoprazole in patients at risk for gastrointestinal bleeding in the ICU. \href{https://doi.org/10.1056/NEJMoa1714919}{\textit{N Engl J Med}. 2018;379:2199-2208.}

    \item Wang Y, Parpia S, Ge L, et al. Proton-pump inhibitors to prevent gastrointestinal bleeding—an updated meta-analysis. \textit{NEJM Evid}. 2024;2:8.

    \item PEPTIC Investigators. Effect of stress ulcer prophylaxis with proton pump inhibitors vs histamine-2 receptor blockers on in-hospital mortality among ICU patients receiving invasive mechanical ventilation: the PEPTIC randomized clinical trial. \href{https://doi.org/10.1001/jama.2019.22190}{\textit{JAMA}. 2020;323:616-626.}
\end{enumerate}

\end{document}
